\scanpage{193}
invariantes par $G$. On y trouve les éléments
\[
 \left.\begin{aligned}
 b_i&=\operatorname{Tr}x_i\\
 c_i&=\det x_i
 \end{aligned}\right\}\qquad1\leq i\leq4
\]
\[
 t_{ij}=\deleted{\unsure{t}}t_{\{i,j\}}=\operatorname{Tr}x_ix_j\qquad\{i,j\}\in\unsure{\mathcal P_2|P_2|\mathfrak P_2}([1,4])
\]
(ce qui fait déjà $4+4+6=14$ fonctions sur un ``espace'' quotient'' de dim $4\times4-3=13$ -- il doit y avoir \deleted{\unsure{donc|deux}} \inserted{une} relation entre ces quantités), et enfin les
\[
 t_{ijk}=\operatorname{Tr}(x_ix_jx_k)\qquad
 \deleted{1\leq i\unsure{<|\leq}j\unsure{<|\leq}k\leq4}\quad
 \inserted{i,j,k\in[1,4]\ \text{tous différents}}
\]
qui sont algébriques sur l'anneau engendré par les $b_i,c_i,t_{ij}$.

Il y a $24$ fonctions $t_{ijk}$, \deleted{mais} (\unsure{correspondants|correspondant} aux \unsure{repères|copies} de l'ens à $4$ éléments) mais
\[
 t_{ijk}=t_{jki}=t_{kij}
\]
donc il n'y a que $\unsure{8|4}$ fonctions distinctes, lesquelles se groupent $2$ par $2$ \deleted{\unsure{en}}
\[
 \left\{\begin{aligned}
 t_{ijk}+t_{jik}&=\Sigma_{\{i,j,k\}}\in\deleted{\unsure{\mathbb Z|k}}\mathbb Z[b_i,c_i,t_i,b_j,c_j,t_j,b_k,c_k,t_k]\\
 t_{ijk}.t_{jik}&=\Pi_{\{i,j,k\}}\in\text{id}
 \end{aligned}\right.\tag*{$(*)$}
\]
% archive sheet number: 192
